\chapter*{Scope and Topic Map}\label{ch:auto-001}
\addcontentsline{toc}{chapter}{Scope and Topic Map}


\section*{How to Use This Text}\label{sec:auto-001}
Chapters 1 through 10 provide topic depth with definitions, equations, and worked examples.
Chapter 5 provides integrated synthesis and cache-simulation context.
Chapter 6 develops dynamic instruction-level parallelism mechanisms such as superscalar scheduling and out-of-order execution.
Chapter 7 develops compiler-scheduled instruction-level parallelism through very long instruction word (VLIW) execution and software pipelining.
Chapter 8 develops multiprocessor organization, coherence, and parallel programming models.
Chapter 9 develops atomic synchronization, transactional-memory ideas, and memory-consistency semantics.
Chapter 10 develops parallel interconnection-network topologies, routing, flow control, and switch fabrics.
Appendix~\ref{ch:advanced-topics} summarizes the advanced-topics lecture on quantum-computing architecture and superconducting processor research.
Appendix~\ref{ch:final-review-checklist} provides a compact final-review checklist and clarification guide.


\section*{Cross-Chapter Reading Paths}\label{sec:cross-chapter-reading-paths}
\begin{enumerate}
    \item Readers who want the performance-and-memory foundation should begin with Chapters 1 through 3 and then read the first part of Chapter 4.
    \item Readers focused on control hazards, prediction, and data hazards should concentrate on the middle and later sections of Chapter 4.
    \item Readers focused on instruction-level parallelism should read Chapter 6 for dynamic ILP and Chapter 7 for compiler-scheduled ILP\@.
    \item Readers focused on simulator-style interpretation and consolidated worked checks should use Chapter 5 alongside the detailed discussions in Chapters 3, 4, and 6.
    \item Readers focused on thread-level parallelism, coherence, and shared-memory multiprocessors should continue from Chapters 1 and 2 into Chapters 8 through 10.
    \item Readers reviewing the advanced-topics final-exam material should read Appendix~\ref{ch:advanced-topics}.
    \item Readers doing a last comprehensive review should use Appendix~\ref{ch:final-review-checklist} to identify weak concepts, then return to the relevant chapter sections.
\end{enumerate}


\section*{Topic Coverage Map}\label{sec:auto-002}

\paragraph{Chapter 1: Introduction}\label{par:auto-001}
Coverage includes speculation with and without prediction, grain sizes of parallelism (bit, instruction, thread, and inter-program), Flynn taxonomy, single instruction, multiple data (SIMD) versus multiple instruction, multiple data (MIMD), pipelining as temporal parallelism, and the distinction between packed-word SIMD and pipelined vector operations.

\paragraph{Chapter 2: Performance Evaluation}\label{par:auto-002}
Coverage includes throughput versus runtime, the central processing unit (CPU)-time equation in terms of instruction count (IC), cycles per instruction (CPI), and cycle time, reduced instruction set computer (RISC) versus complex instruction set computer (CISC) tradeoffs, Amdahl's and Gustafson's laws, and correct use of arithmetic versus harmonic means.

\paragraph{Chapter 3: Cache Memories}\label{par:auto-003}
Coverage includes static random-access memory (SRAM) versus dynamic random-access memory (DRAM) and bit-line delay, locality, hierarchy levels, virtual-memory terminology and its analogy to caching, average access time (AAT), the three C's of misses, cache geometry and construction, replacement and write policies, insertion-policy choices such as MIP and LIP, prefetching strategies including Markov and hybrid schemes, techniques to improve miss rate, miss penalty, and hit time, simulator-style hierarchy statistics, and compiler/data-layout effects.

\paragraph{Chapter 4: Pipelining}\label{par:auto-004}
Coverage includes instruction fetch / instruction decode / execute / memory / write back (IF/ID/EX/MEM/WB) behavior, speedup equations, cycles per instruction (CPI) stall modeling, the principle that program behavior plus hazards produces stalls, structural and control hazards, delay and squash-slot ideas, branch target buffer (BTB) behavior with and without tags, return-address stack (RAS) behavior, direction-prediction heuristics, and dynamic predictors including one-bit and Smith-counter schemes, GShare, GSelect, Yeh/Patt, perceptrons, hybrid tournament selection, and delayed predictor-update timing.

\paragraph{Chapter 5: Integrated Synthesis and Cache Simulation Concepts}\label{par:auto-005}
Coverage includes integrated synthesis, worked pipeline calculations, hierarchy semantics, prefetching, and interpretation of simulator-style cache, branch, and average access time (AAT) statistics.

\paragraph{Chapter 6: Dynamic ILP, Superscalar Scheduling, and Out-of-Order Execution}\label{par:scope-ch6}
Coverage includes cycles per instruction (CPI) versus instructions per cycle (IPC) interpretation, instruction-level parallelism (ILP), historical superscalar context, fire/complete ordering models, scoreboarding and busy-bit scheduling, function-unit latency considerations, Tomasulo-style reservation stations, physical register files, result buses, register renaming intuition, precise versus imprecise exceptions, reorder buffer (ROB) and future-file organization, checkpoint/repair recovery, simulator-style fetched/retired and queue-pressure counter interpretation, and the security consequences of speculative out-of-order execution.

\paragraph{Chapter 7: Compiler-Scheduled ILP and Software Pipelining}\label{par:scope-ch7}
Coverage includes compiler-scheduled very long instruction word (VLIW) execution, dependence-graph and list-scheduling ideas, zero-no-op multi-op encoding with tail and pause fields, software pipelining and iterative modulo scheduling, and memory disambiguation.

\paragraph{Chapter 8: Multiprocessors, Interconnects, and Cache Coherence}\label{par:scope-ch8}
Coverage includes shared-memory versus message-passing organization, standalone multiprocessors versus accelerator systems, buses versus point-to-point links, the cache-coherence problem and its formal properties, shared-cache versus directory-based versus snooping solutions, MSI coherence states and transactions, MESI/MOESI/MOESIF refinements including \texttt{E}, \texttt{O}, and \texttt{F}, transient coherence states, directory-entry structure and message flow, limited and distributed directory organizations, sharing patterns and false sharing, and shared-memory programming support through locks, barriers, and atomic synchronization primitives.

\paragraph{Chapter 9: Atomic Synchronization, Transactional Memory, and Memory Consistency}\label{par:scope-ch9}
Coverage includes the distinction between coherence states and full protocol behavior, abort-versus-intervene transfer policies, load-link / store-conditional, fetch-and-op and fetch-and-add primitives, dynamic work claiming, tree and tournament locks, sense-reversing barriers, hardware barrier intuition, transactional-memory execution and aborts, and the distinction between coherence and memory consistency, including strict consistency, sequential consistency, relaxed consistency, data races, and synchronization ordering through fences and acquire/release style semantics.

\paragraph{Chapter 10: Parallel Interconnection Networks}\label{par:scope-ch10}
Coverage includes the distinction between direct and indirect topologies, hypercube organization and Gray-code intuition, $k$-ary $n$-dimensional meshes and tori, physical-layout constraints and wire-length balance, bisection bandwidth, crossbars and dance-hall style switching, multi-stage interconnects such as omega networks, perfect-shuffle wiring, deterministic versus adaptive versus oblivious routing, packets versus flits, and the progression from store-and-forward to cut-through, wormhole, and virtual-channel flow control.

\paragraph{Appendix D: Advanced Topics}\label{par:scope-appd}
Coverage includes quantum computers as accelerators, decoherence, imperfect control, NISQ limits, fault-tolerant logical qubits, reversible phase-oracle intuition, superconducting-circuit constraints, compiler-scheduled dataflow execution, operand queues, barrel processing, and combined single-thread/throughput-oriented superconducting processor organization.

\paragraph{Appendix E: Final Review Checklist}\label{par:scope-appe}
Coverage includes a comprehensive concept checklist, review clarifications for tiling, iterative modulo scheduling, memory consistency, barriers versus locks, and bisectional bandwidth, plus fast self-check prompts.
